\setupcolors[state=start]
\setupcolor[hex]
\definecolor[darkblue][h=003366]
\definecolor[darkgray][h=e6e6e6]
\definecolor[redline][h=d9534f]
\definecolor[blueline][h=5bc0de]
\definecolor[greenline][h=5cb85c]

\setuphead[section][style=\bfa, color=darkblue, before=\blank, after=\blank]
\setuphead[subsection][style=\bf, color=black, before=\blank, after=\blank]
\setuppapersize[A4][A4]
\setupbodyfont[11pt]
\setuplayout[width=middle, height=middle, topspace=2cm, bottomspace=2cm]

\starttext

\title{\bf Avances: heurísticas y modelo exacto en dos etapas}
\blank[big]

Este documento resume el estado de los algoritmos implementados para el
problema de picking y ruteo en almacén (OPP): dos heurísticos (Hill Climbing y Ant Colony Optimization),
cada uno con una variante base y una variante que incorpora congestión
entre pickers concurrentes, y \type{opp_v5}, una reformulación en dos
etapas del modelo exacto que agrega velocidad heterogénea de pickers y un
límite de cadena de frío. Los cuatro heurísticos resuelven exactamente el
mismo problema que el modelo exacto (\type{opp_v4}), con las mismas
restricciones, verificado por un validador explícito en cada uno.

\startsection[title={Hill Climbing (HC)}]

Búsqueda local: parte de una solución inicial factible y explora vecinos
mediante movimientos que preservan factibilidad, aceptando solo mejoras
(sin annealing ni tabú). Es determinístico salvo por la elección aleatoria
de \emph{qué} movimiento probar en cada paso. No hay ningún mecanismo
probabilístico en la evaluación en sí.

\startsubsection[title={Algoritmo base (\type{opp\_hc\_2})}]

\starttyping[option=none]
ALGORITMO Solve(params)
    S <- GenerarSolucionInicialFactible()
    si S es None: retornar infeasible

    costo <- Evaluar(S)
    iteracion <- 0 ; estancamiento <- 0

    mientras iteracion < Max_it y estancamiento < Max_imp
                                y tiempo_transcurrido < Max_time:
        iteracion <- iteracion + 1
        S' <- GenerarVecino(S)
        si S' es None:
            estancamiento <- estancamiento + 1 ; continuar

        costo' <- Evaluar(S')
        si costo' < costo:                    # steepest descent, sin annealing
            S, costo <- S', costo'
            estancamiento <- 0
        si_no:
            estancamiento <- estancamiento + 1

    retornar S, costo


SUBRUTINA GenerarSolucionInicialFactible()
    S <- solución vacía
    para cada lote en B con órdenes:
        ruta <- DijkstraVecinoMasCercano(ubicaciones_requeridas(lote))
        si ruta es None:
            ruta <- BusquedaExhaustivaBFS(ubicaciones_requeridas(lote))
        si ruta es None:
            retornar None                     # grafo desconectado / inalcanzable
        S.fijar_ruta(lote, ruta)

    asignar cada lote con órdenes a un picker, round-robin (X_pb)
    retornar S si Validar(S) si_no None


SUBRUTINA GenerarVecino(S)
    elegir al azar entre los movimientos aplicables:

    2-opt en una ruta:
        lote <- lote al azar con ruta no vacía
        mejor <- None ; mejor_distancia <- distancia(ruta(lote))
        para cada segmento [i,j] a invertir en la ruta:
            candidata <- ruta con [i,j] invertido
            si candidata es factible (requeridas cubiertas, arcos válidos,
                                       sin repetir nodos)
               y distancia(candidata) < mejor_distancia:
                mejor, mejor_distancia <- candidata, distancia(candidata)
        retornar S con ruta(lote) <- mejor  (o None si ninguna mejora)

    Reasignar picker:
        lote <- lote al azar con órdenes
        picker <- picker al azar distinto del actual
        retornar S con lote reasignado (rutas sin tocar)

    Intercambiar secuencia:
        picker <- picker al azar con >= 2 lotes en cola
        intercambiar el orden de dos lotes al azar en esa cola
        retornar S con esa cola reordenada
\stoptyping

La construcción de rutas (\type{DijkstraVecinoMasCercano}) es \emph{greedy}
y no tiene mecanismo de retroceso: en instancias con muchas ubicaciones
requeridas por lote puede quedar atrapada en un callejón sin salida del
que no puede deshacer una elección anterior. El respaldo exhaustivo
(\type{BusquedaExhaustivaBFS}) explora en dos niveles: una búsqueda
rápida memoizada primero, y solo si falla, una búsqueda sin poda como
último recurso, acotada.

\stopsubsection

\startsubsection[title={Qué agrega \type{opp\_hc\_3}: congestión}]

La construcción de rutas y los tres movimientos de vecindad son
idénticos a \type{opp\_hc\_2}: la congestión entra únicamente en
\type{Evaluar}, no en cómo se construyen o modifican las rutas:

\starttyping[option=none]
SUBRUTINA Evaluar(S)                    # reemplaza a Evaluar en hc_2
    tracker <- nuevo CongestionTracker(Min_dwell)
    tiempo_picker[p] <- 0  para cada picker p

    # Round-robin por SLOT de cola de picker: un lote se precia contra todo
    # lo que ya reservaron los lotes que, dadas las colas fijas, corren
    # antes o al mismo tiempo (misma aproximación causal que opp_aco_3).
    para cada slot en 0 .. max(largo de cola sobre todos los pickers) - 1:
        para cada picker en P (orden determinístico):
            si el slot no tiene lote para este picker: continuar
            lote <- S.secuencia[picker][slot]
            ruta <- S.ruta(lote)

            base <- B_setup + B_picking*capacidad(lote)
                    + B_time*distancia(ruta) + tiempo_servicio(ruta)
            espera <- SimularCongestionDeRuta(ruta, dueño=picker,
                                               inicio=tiempo_picker[picker],
                                               tracker)

            tiempo_picker[picker] <- tiempo_picker[picker] + base + espera
            completitud[lote] <- tiempo_picker[picker]

    adelanto, tardanza <- comparar completitud[lote] vs. Dd(orden)
    retornar W_full*sum(completitud) + W_ea_ta*(W_ea*adelanto + W_ta*tardanza)
\stoptyping

\type{CongestionTracker} modela cada nodo y cada arco como un recurso de
capacidad 1 (adaptado de Chen et al., 2013): reservar un recurso ya
ocupado hace esperar al picker que llega después. Como la ruta se
construye igual que en \type{opp\_hc\_2} y la congestión solo se \emph{precia}
al final, el costo de agregar esta capacidad es bajo.

\stopsubsection

\stopsection

\startsection[title={Ant Colony Optimization (ACO)}]

Metaheurística poblacional: en cada iteración, varias ''hormigas''
construyen una solución completa cada una, guiadas por feromona
depositada por hormigas anteriores más una heurística de visibilidad
(distancia). Al final de la iteración, la feromona se evapora y se
refuerza según qué tan buenas fueron las soluciones encontradas.

\startsubsection[title={Algoritmo base (\type{opp\_aco\_2})}]

\starttyping[option=none]
ALGORITMO Solve(params)
    tau       <- inicializar feromona de arco (todos los arcos válidos)
    tau_picker <- inicializar feromona (lote, picker) en 1.0
    mejor_solucion <- None ; mejor_costo <- +inf

    para iteracion en 1 .. Max_it:
        soluciones_hormiga <- []

        para hormiga en 1 .. A_ants:
            S <- ConstruirSolucionHormiga(tau, tau_picker)
            si S no es None y Validar(S):
                soluciones_hormiga.agregar( (S, Evaluar(S)) )

        si soluciones_hormiga está vacío: continuar

        mejor_iteracion <- el (S, costo) con menor costo
        si mejor_iteracion.costo < mejor_costo:
            mejor_solucion, mejor_costo <- copia(mejor_iteracion)

        ActualizarFeromonas(tau, tau_picker, soluciones_hormiga,
                             mejor_solucion, mejor_costo)

    retornar mejor_solucion, mejor_costo


SUBRUTINA ConstruirSolucionHormiga(tau, tau_picker)
    S <- solución vacía
    carga_picker[p] <- 0 para cada picker p

    para cada lote en B con órdenes (orden fijo):
        picker <- SeleccionRuleta(lote, tau_picker, carga_picker)  # favorece
        S.asignar(lote, picker)                                    # pickers
        carga_picker[picker] += 1                                  # con menos carga

        ruta <- None
        repetir hasta Max_route_attempts_per_batch veces:
            candidata <- ConstruirRuta(ubicaciones_requeridas(lote), tau)
            si candidata no es None y candidata no tiene loops:
                ruta <- candidata ; salir del repetir
        si ruta es None:
            retornar None          # la hormiga entera se descarta
        S.fijar_ruta(lote, ruta)

    retornar S


SUBRUTINA ConstruirRuta(requeridas, tau)
    ruta <- [deposito] ; visitados <- {deposito} ; actual <- deposito
    faltantes <- copia(requeridas)

    mientras faltantes no esté vacío:
        vecinos <- vecinos(actual) no visitados
        si vecinos está vacío:               # respaldo: permitir revisitar
            vecinos <- vecinos(actual) sin filtrar
            si vecinos está vacío: retornar None

        para cada vecino n:
            tau_n <- tau(actual, n) ^ Alpha
            si n en faltantes:
                eta_n <- (1.5 / distancia(actual, n)) ^ Beta   # prioridad alta
            si_no:
                eta_n <- (1 / (distancia(actual,n)
                              + distancia_a_mas_cercana_faltante(n))) ^ Beta
            peso(n) <- tau_n * eta_n

        siguiente <- SeleccionRuleta(vecinos, pesos)
        ruta.agregar(siguiente) ; visitados.agregar(siguiente)
        faltantes.quitar(siguiente) ; actual <- siguiente

    ruta.agregar( camino_mas_corto(actual -> final) )   # BFS determinístico
    retornar ruta


SUBRUTINA ActualizarFeromonas(tau, tau_picker, soluciones, mejor, costo_mejor)
    evaporar tau y tau_picker por factor (1 - E)
    para (S, costo) en soluciones:
        depositar Q/costo en tau, para cada arco que S usó
        depositar Q/(1+costo) en tau_picker[(lote,picker)], por cada lote
    si mejor existe:                                   # refuerzo elitista
        depositar boost extra en tau y en tau_picker para la mejor solución
\stoptyping

La construcción de rutas usa un presupuesto acotado de reintentos por
lote (\type{Max\_route\_attempts\_per\_batch}, 100 por defecto) con
revisita como último recurso. Sin esto, exigir que \emph{todos} los lotes de una hormiga
salieran sin loops simultáneamente volvía la construcción
infeasible a partir de cierta cantidad de lotes.

\stopsubsection

\startsubsection[title={Qué agrega \type{opp\_aco\_3}: congestión y arquitectura multi-feromona}]

A diferencia de \type{opp\_hc\_3}, acá la congestión no se agrega solo al
evaluar: se incorpora \emph{dentro} de la construcción de la ruta, paso a
paso, con una segunda tabla de feromona dedicada (arquitectura m-ACO4(1,m)
de Alaya et al., con la parametrización calibrada de Lesch et al., 2023):

\starttyping[option=none]
SUBRUTINA ConstruirRuta_con_congestion(requeridas, tau1, tau2, tracker,
                                        dueño, tiempo_inicio)
    ruta <- [deposito] ; actual <- deposito ; t <- tiempo_inicio
    faltantes <- copia(requeridas)

    mientras faltantes no esté vacío:
        usar_tau1 <- ( azar() < P_objective_switch )   # UN objetivo por paso,
                                                          # compartido por todos
                                                          # los candidatos
        para cada vecino n de actual, no visitado:
            cercania(n) <- heurística de distancia hacia lo faltante
            espera_estimada <- tracker.espera_estimada(arco(actual,n)
                                                         [y nodo(n)], en t)
            si usar_tau1:
                peso(n) <- tau1(actual,n)^Alpha * cercania(n)^Beta
            si_no:
                # cercania(n) también entra acá (fix real, ver más abajo)
                peso(n) <- tau2(actual,n)^Alpha *
                           ( cercania(n) * 1/(1+espera_estimada) )^Beta2

        siguiente <- SeleccionRuleta(vecinos, pesos)
        t, espera <- tracker.reservar(arco(actual,siguiente), nodo(siguiente),
                                       en t, dueño)
        ruta.agregar(siguiente) ; actual <- siguiente

    ruta.agregar( camino_mas_corto(actual -> final) )
    retornar ruta, espera_total
\stoptyping

Además: el orden de procesamiento de los lotes de un picker se ordena por
fecha de vencimiento más próxima antes de construir rutas (en vez de
quedar en orden arbitrario de asignación); hay un mecanismo de
''cataclismo'' que reinicia ambas tablas de feromona tras varias
iteraciones sin mejora (Chen et al., 2013; Lesch et al., 2023); y las
feromonas están acotadas en un rango \type{[Tau\_min, Tau\_max]} (MAX-MIN
Ant System).

\blank

\stopsubsection

\stopsection

\startsection[title={\type{opp\_v5}: modelo exacto en dos etapas}]

\type{opp\_v5} reformula el modelo exacto monolítico de \type{opp\_v4} en
dos etapas resueltas en secuencia con SCIP, agregando dos capacidades que
\type{opp\_v4} no tiene: velocidad heterogénea por picker y un límite de
cadena de frío. La descomposición es válida porque el ruteo no depende de
qué picker ejecuta el lote (no hay congestión modelada), así que
optimizar cada lote de forma independiente en la Etapa 1 es equivalente a
optimizarlo dentro del problema conjunto.

\startsubsection[title={Etapa 1: ruteo por lote (independiente)}]

Para cada lote, resuelve el camino de menor tiempo (traslado + servicio)
que visita sus ubicaciones requeridas, sin considerar todavía qué
picker lo va a ejecutar. Usa la misma familia de restricciones que
\type{opp\_v4} para garantizar una ruta simple (salida del depósito,
llegada al nodo final, conservación de flujo, sin ciclos), con
eliminación de subtours vía una formulación de flujo de un solo producto
(Gavish-Graves) en vez de la que usa \type{opp\_v4}. Exporta, para cada
lote $b$, el tiempo óptimo de ruta $Rt_b$.

\stopsubsection

\startsubsection[title={Etapa 2: asignación, secuenciación y cadena de frío}]

Toma $Rt_b$ como dato fijo (ya no es variable de decisión) y decide qué
picker ejecuta cada lote, en qué orden, respetando adelanto/tardanza. El
tiempo de ejecución de un lote en un picker,
$t\_exec_{pb} = F_b + Rt_b / V_p$ (donde $F_b$ es el tiempo fijo de
\emph{setup}+\emph{picking}, independiente del picker, y $V_p$ es el
factor de velocidad del picker $p$), pasa a ser un parámetro
derivado en vez de una variable (ya no hace falta linealizar con Big-M
para vincularlo). La cadena de frío se implementa restringiendo el
dominio de las variables de asignación al conjunto de pares
elegibles $\mathcal{PB}_{elig} = \{(p,b) : t\_exec_{pb} \le T_{frio}\}$, en
vez de agregar una restricción aparte.

\blank

La definición formal completa (conjuntos, parámetros, variables, función
objetivo y las 9 restricciones de la Etapa 2) está en
\type{opp\_v5.tex}/\type{.pdf}, con el mismo formato usado para
\type{opp\_model\_v4.tex}. Función objetivo de la Etapa 2 (idéntica en
forma a la de \type{opp\_v4}, aplicada sobre $v_b$ en vez de sobre
variables de ruteo):

\placeformula
\startformula
\min \left( W_{full} \cdot \sum_{b \in \mathcal{B}} v_{b} + W_{ea\_ta} \cdot \left( W_{ea} \cdot \sum_{n \in \mathcal{N}} ea_{n} + W_{ta} \cdot \sum_{n \in \mathcal{N}} ta_{n} \right) \right)
\stopformula

\stopsubsection

\startsubsection[title={Estado y validación}]

\type{opp\_v5} todavía está en desarrollo. Sobre instancias de test unitarias (es decir, no de estrés), sus
19 casos comparables contra \type{opp\_v4} dan un gap de calidad
indistinguible de cero (ruido de punto flotante). Esta es una validación cruzada de
que la descomposición en dos etapas no pierde optimalidad frente a la
formulación monolítica, en el régimen evaluado hasta ahora. Se agregaron tests que validan específicamente sus dos
capacidades nuevas.

Su batería de estrés, ahora completa (8 instancias), da 0\% de éxito
(7 timeouts y 1 falla real). Por lo tanto, no es competitivo a esa escala frente al
12.5\% de \type{opp\_v4}, así que se mantiene fuera de las comparaciones
de escalabilidad. La falla real (no por timeout) reveló un segundo punto
a revisar: en al menos una instancia grande,
la Etapa 1 reporta una ruta óptima con muchos más arcos activos de los
necesarios y ninguno llega al nodo final, pese a que el propio solver la
marca como óptima.

\stopsubsection

\stopsection

\stoptext
